% Part: computability
% Chapter: computability-theory
% Section: rice-theorem

\documentclass[../../../include/open-logic-section]{subfiles}

\begin{document}

\olfileid[hi]{cmp}{thy}{rce}
\olsection{राइस का प्रमेय}

विचार करने पर आप देखेंगे कि $\fn{Tot}$ की विशिष्टताओं की
\olref[tot]{prop:total} के प्रमाण में कोई भूमिका नहीं है। हमने $h(x,y)$
को इस प्रकार बनाया था कि वह अचर फलन $j(y) = 0$ की तरह ठीक तब व्यवहार
करे जब $x$, $K$ में हो; किंतु उन परिस्थितियों में हम उससे किसी
अन्य आंशिक संगणनीय फलन जैसा व्यवहार भी करा सकते थे। यह अवलोकन हमें एक
अधिक सामान्य प्रमेय कहने देता है, जिसका मोटा आशय है कि संगणनीय फलनों का
कोई अतुच्छ गुण निर्णेय नहीं है।

ध्यान रखें कि $\cfind{0}$, $\cfind{1}$, $\cfind{2}$,~\dots आंशिक
संगणनीय फलनों का हमारा मानक परिगणन है।

\begin{thm}[राइस का प्रमेय]
  $C$ को आंशिक संगणनीय फलनों का कोई समुच्चय मानें और $A =
  \Setabs{n}{\cfind{n} \in C}$ मानें। यदि $A$ संगणनीय है, तो या तो $C$
  रिक्त है अथवा $C$ सभी आंशिक संगणनीय फलनों का समुच्चय है।
\end{thm}

{\em सूचकांक समुच्चय} वह समुच्चय $A$ है जिसमें यह गुण होता है कि यदि
$n$ और $m$ ऐसे सूचकांक हैं जो समान फलन को ``संगणित'' करते हैं, तो या तो
$n$ और $m$ दोनों $A$ में हैं, अथवा दोनों में से कोई नहीं। यह देखना कठिन
नहीं कि प्रमेय के समुच्चय~$A$ में यह गुण है। विलोमतः, यदि $A$~सूचकांक
समुच्चय और $C$~इन सूचकांकों द्वारा संगणित फलनों का समुच्चय है, तो
$A = \Setabs{n}{\cfind{n} \in C}$।

\begin{explain}
इस शब्दावली में राइस का प्रमेय यह कहने के समतुल्य है कि कोई अतुच्छ
सूचकांक समुच्चय निर्णेय नहीं है। प्रमेय का अर्थ समझने के लिए
\emph{प्रोग्रामों} (मान लें, आपकी पसंदीदा प्रोग्रामिंग भाषा में) और उनके
द्वारा संगणित फलनों का भेद उभारना सहायक है। प्रोग्राम (सूचकांक) वाक्यगत
वस्तुएँ हैं, और उनके बारे में कुछ प्रश्न निश्चय ही संगणनीय हैं: क्या इस
प्रोग्राम में 150 से अधिक चिह्न हैं? क्या इसमें 22 से अधिक पंक्तियाँ हैं?
क्या इसमें ``while'' कथन है? क्या ``hello world'' अक्षर-शृंखला कभी किसी
``print'' कथन के तर्क के रूप में आती है? राइस का प्रमेय कहता है कि
प्रोग्राम के \emph{व्यवहार} के बारे में कोई अतुच्छ प्रश्न संगणनीय नहीं
है। इनमें ऐसे प्रश्न आते हैं: क्या प्रोग्राम आगत $0$ पर रुकता है? क्या
वह कभी रुकता है? क्या वह कभी कोई सम संख्या निर्गत करता है?
\end{explain}

\begin{proof}[राइस के प्रमेय का प्रमाण]
मान लें कि $C$ न तो रिक्त है, न सभी आंशिक संगणनीय फलनों का समुच्चय, और
$A$ को~$C$ के फलनों के सूचकांकों का समुच्चय मानें। हम दिखाएँगे कि यदि
$A$ संगणनीय होता, तो हम विराम समस्या हल कर सकते थे; अतः $A$~संगणनीय नहीं है।

व्यापकता की हानि के बिना हम मान सकते हैं कि वह फलन $f$ जो कहीं भी
परिभाषित नहीं है, $C$ में नहीं है (अन्यथा नीचे के तर्क में $C$ और उसके
पूरक की अदला-बदली करें)। $g$ को~$C$ का कोई फलन मानें। विचार यह है कि यदि
हम~$A$ का निर्णय कर सकें, तो~$f$ को संगणित करने वाले सूचकांकों और~$g$ को
संगणित करने वाले सूचकांकों का भेद बता सकेंगे; फिर इस क्षमता से विराम
समस्या हल कर सकेंगे।

इस प्रकार। सार्वत्रिक आंशिक संगणनीय फलनों का उपयोग करके हम एक फलन
परिभाषित कर सकते हैं:
\[
h(x,y) \simeq
\begin{cases}
\text{अपरिभाषित} & \text{यदि $\cfind{x}(x) \fundefined$} \\
g(y) & \text{अन्यथा।}
\end{cases}
\]
$h$ को संगणित करने के लिए पहले हम $\cfind{x}(x)$ को संगणित करने का
प्रयत्न करते हैं; यदि वह संगणना रुकती है, तो~$g(y)$ को संगणित करते हैं;
और यदि {\em वह} संगणना रुकती है, तो निर्गत लौटाते हैं। अधिक औपचारिक रूप
में हम लिख सकते हैं
\[
h(x,y) \simeq \Proj{2}{0}(g(y),\fn{Un}(x,x)).
\]
जहाँ $\Proj{2}{0}(z_0, z_1) = z_0$ वह $2$-स्थानी प्रक्षेपण फलन है जो
$0$-वाँ तर्क लौटाता है; यह संगणनीय है।

तब $h$ आंशिक संगणनीय फलनों का संयोजन है, और दायाँ पक्ष ठीक तभी परिभाषित
और~$g(y)$ के बराबर है जब $\fn{Un}(x,x)$ तथा $g(y)$ दोनों परिभाषित हों।

ध्यान दें कि किसी नियत~$x$ के लिए, यदि $\cfind{x}(x)$ अपरिभाषित है, तो
$h(x,y)$ प्रत्येक~$y$ पर अपरिभाषित है; और यदि $\cfind{x}(x)$ परिभाषित
है, तो $h(x,y) \simeq g(y)$। अतः~$x$ के किसी नियत मान पर $h(x,y)$ या तो
ठीक $f$ जैसा, या ठीक $g$ जैसा व्यवहार करता है; और $\cfind{x}(x)$ के
परिभाषित होने का निर्णय करना इन्हीं दो स्थितियों में से एक का निर्णय
करना है। किंतु यह निर्णय करना है कि $h_x(y) \simeq h(x,y)$,~$C$ में है
या नहीं; और यदि $A$ संगणनीय होता, तो हम ठीक यही कर सकते थे।

अधिक औपचारिक रूप में, चूँकि $h$~आंशिक संगणनीय है, इसलिए वह फलन
$\cfind{e}$ के बराबर है, किसी सूचकांक~$e$ के लिए। $s$-$m$-$n$ प्रमेय से कोई आदिम
पुनरावर्ती फलन~$s$ ऐसा है कि प्रत्येक~$x$ पर
$\cfind{s(e,x)}(y) = h_x(y)$। अब प्रत्येक $x$ के लिए, यदि
$\cfind{x}(x) \fdefined$, तो $\cfind{s(e,x)}$,~$g$ के समान फलन है और
इसलिए $s(e,x)$, $A$ में है। दूसरी ओर, यदि $\cfind{x}(x)
\uparrow$, तो
$\cfind{s(e,x)}$,~$f$ के समान फलन है और इसलिए $s(e,x)$,~$A$ में नहीं है।
दूसरे शब्दों में, प्रत्येक~$x$ के लिए $x
\in K$ तभी और केवल तभी जब
$s(e,x) \in A$। यदि $A$ संगणनीय होता, तो $K$~भी होता—यह विरोधाभास है।
अतः $A$ संगणनीय नहीं है।
\end{proof}

राइस का प्रमेय अत्यंत शक्तिशाली है। अगला तात्कालिक उपप्रमेय इसके कुछ
नमूना अनुप्रयोग दिखाता है।
\begin{cor}
निम्नलिखित समुच्चय अनिर्णेय हैं।
\begin{enumerate}
\item $\Setabs{x}{\text{$17$, $\cfind{x}$ के परिसर में है}}$
\item $\Setabs{x}{\text{$\cfind{x}$ अचर है}}$
\item $\Setabs{x}{\text{$\cfind{x}$ पूर्ण है}}$
\item $\Setabs{x}{\text{जब भी $y < y'$, $\cfind{x}(y) \fdefined$, और
    यदि $\cfind{x}(y') \fdefined$, तो $\cfind{x}(y) < \cfind{x}(y')$}}$
\end{enumerate}
\end{cor}

\begin{proof}
  ये सभी अतुच्छ सूचकांक समुच्चय हैं।
\end{proof}

\end{document}

